\section*{Problem 4: Critical-Word First and Early Restart (16 points)}

\paragraph{Common setup and notation.}
A word is 8~Bytes, so a block of $B$ Bytes holds $n=B/8$ words. For a given cache
level write
\begin{description}\itemsep2pt
  \item[$t_1$] time until the \emph{first} word arrives when the block is sent in
        normal order (no critical-word first);
  \item[$t_x$] time for each \emph{additional} word;
  \item[$t_c$] time until the requested word arrives when critical-word first is
        used (it is sent first, so this single transfer is all the processor waits for).
\end{description}
Without reordering, word $i$ of the block ($i=0,1,\dots,n-1$) is available after
$t_1+i\,t_x$. The three schemes are then
\begin{align}
  \text{neither:}\quad & T_{\text{none}} = t_1+(n-1)\,t_x
    &&\text{(wait for the whole block)} \label{p4:eq:none}\\
  \text{early restart only:}\quad & T_{\text{ER}} = \sum_{i=0}^{n-1} p_i\,(t_1+i\,t_x)
    &&\text{(resume when word $i$ arrives)} \label{p4:eq:er}\\
  \text{both:}\quad & T_{\text{both}} = t_c
    &&\text{(requested word first, then resume)} \label{p4:eq:both}
\end{align}
where $p_i$ is the probability that the referenced word sits at position $i$.
Equations~\eqref{p4:eq:none} and~\eqref{p4:eq:both} do not involve $p_i$: in the first
the processor waits for the last word regardless of what it asked for, and in the third
the requested word is always the first one delivered. Only~\eqref{p4:eq:er} is a genuine
expectation over the reference distribution.

\subsection*{Question A (8 points)}

Here $B=64$~Bytes, so
\[
  n=\frac{64\ \text{Bytes}}{8\ \text{Bytes/word}}=8\ \text{words per block},
\]
with $t_1=90$, $t_x=10$ and $t_c=110$ CCs. (The 1~MB capacity is not needed: the miss
service time depends on the block size, not on how many blocks the cache holds.) Word
$i$ therefore arrives at $90+10i$ CCs, i.e.\ at
$90,100,110,120,130,140,150,160$ CCs for $i=0,\dots,7$.

The reference distribution is \emph{not} uniform. The first word carries $30\%$ and the
remaining $n-1=7$ words share the other $70\%$ equally:
\[
  p_0=0.30,\qquad
  p_i=\frac{0.70}{7}=0.10\quad(i=1,\dots,7),
  \qquad\text{check: } 0.30+7(0.10)=1 .
\]

\paragraph{(1) Neither critical-word first nor early restart.}
The processor resumes only once the entire block has arrived, which happens when the
last word lands:
\[
  T_{\text{none}}=t_1+(n-1)t_x=90+7\times 10=\boxed{160\ \text{CCs}} .
\]
This does not depend on which word was requested, so the reference distribution is
irrelevant here and the ``average'' equals this fixed value.

\paragraph{(2) Early restart only.}
The block still arrives in normal order, but the processor resumes as soon as the
requested word has arrived, so the wait depends on that word's position:
\begin{align*}
  T_{\text{ER}}
  &= 0.30\times 90
   + 0.10\bigl(100+110+120+130+140+150+160\bigr)\\
  &= 27 + 0.10\times 910
   = 27+91
   = \boxed{118\ \text{CCs}} .
\end{align*}

\paragraph{(3) Both critical-word first and early restart.}
Memory delivers the requested word first, at a cost of $t_c=110$ CCs, and the processor
restarts immediately:
\[
  T_{\text{both}}=110\ \text{CCs}\ \Rightarrow\ \boxed{110\ \text{CCs}} .
\]
Again the block-position distribution drops out entirely: whichever word is requested,
it is the one sent first, so every miss costs the same 110 CCs. Note that
critical-word first is \emph{more expensive per transfer} than a normal first word
($110>90$); it still wins overall because it removes the wait for all seven remaining
words.

\paragraph{Summary.}
\begin{center}
\begin{tabular}{lcc}
\toprule
Scheme & Average L2 miss service time & Saving vs.\ (1)\\
\midrule
(1) Neither                       & 160 CCs & ---\\
(2) Early restart only            & 118 CCs & $42/160=26.25\%$\\
(3) Critical-word first + early restart & 110 CCs & $50/160=31.25\%$\\
\bottomrule
\end{tabular}
\end{center}

\subsection*{Question B (4 points)}

\textbf{Answers: the first blank is $18$ CCs and the second blank is $10\%$.}

Now $B=32$~Bytes, so
\[
  n=\frac{32\ \text{Bytes}}{8\ \text{Bytes/word}}=4\ \text{words per block},
\]
with $t_x=4$ CCs, $t_c=27$ CCs, and every word equally likely, $p_i=\tfrac14$. Let $x$
denote the unknown first-word latency.

\paragraph{First blank --- recovered from the stated $20\%$.}
The paragraph gives us the early-restart result and asks us to work backwards. In terms
of $x$,
\begin{align*}
  T_{\text{none}} &= x+(n-1)t_x = x+3\times 4 = x+12,\\
  T_{\text{ER}}   &= \frac14\bigl[x+(x+4)+(x+8)+(x+12)\bigr]
                   = x+\frac{4(0+1+2+3)}{4} = x+6 .
\end{align*}
Early restart alone reduces the original time by $20\%$, so
\[
  \frac{T_{\text{none}}-T_{\text{ER}}}{T_{\text{none}}}
  =\frac{(x+12)-(x+6)}{x+12}
  =\frac{6}{x+12}=0.20
  \ \Longrightarrow\ x+12=30
  \ \Longrightarrow\ \boxed{x=18\ \text{CCs}} .
\]
Hence $T_{\text{none}}=30$ CCs and $T_{\text{ER}}=24$ CCs, which indeed differ by
$6/30=20\%$.

\paragraph{Second blank.}
With critical-word first the requested word arrives after $t_c=27$ CCs and the
processor restarts there, so $T_{\text{both}}=27$ CCs and
\[
  \frac{T_{\text{none}}-T_{\text{both}}}{T_{\text{none}}}
  =\frac{30-27}{30}=\frac{1}{10}=\boxed{10\%} .
\]
Both blanks are exact integers, as the question requires.

\subsection*{Question C (4 points)}

\paragraph{Verdict.} On these numbers the two techniques are \textbf{more important for
the L2 cache}.

\paragraph{Reason (from the results of A and B).}
Comparing the savings the two techniques actually deliver at each level:
\begin{center}
\begin{tabular}{lcc}
\toprule
 & L2 (Question A) & L1 (Question B)\\
\midrule
Words per block                     & 8       & 4\\
Neither                             & 160 CCs & 30 CCs\\
Early restart only                  & 118 CCs \;($-26.25\%$) & 24 CCs \;($-20\%$)\\
Critical-word first + early restart & 110 CCs \;($-31.25\%$) & 27 CCs \;($-10\%$)\\
Absolute saving of the pair         & 50 CCs  & 3 CCs\\
\bottomrule
\end{tabular}
\end{center}
At L2 the pair of techniques removes $31.25\%$ of the miss service time (50 CCs); at L1
it removes only $10\%$ (3 CCs). The gap is even sharper for critical-word first taken on
its own: at L2 it improves on early restart alone ($118\to110$ CCs), whereas at L1 it
makes matters \emph{worse} ($24\to 27$ CCs), because its $27-18=9$~CC surcharge exceeds
the 6~CC that early restart had already saved. So at L1 it would be better not to enable
critical-word first at all, while at L2 it pays for itself --- which is precisely what
``more important for L2'' means.

\paragraph{Factors that drive the relative importance.}
\begin{enumerate}\itemsep3pt
  \item \emph{Words per block, $n$.} The most these techniques can ever save is the
        transfer tail $(n-1)t_x$ --- the time spent shipping the words the processor is
        not waiting for. Larger blocks mean a longer tail and more to skip. Lower cache
        levels normally use larger blocks (8 words here versus 4), which favours L2.
  \item \emph{Fixed first-word latency versus the tail.} The ceiling on the saving is
        $(n-1)t_x/\bigl[t_1+(n-1)t_x\bigr]$: $70/160=43.75\%$ at L2 and $12/30=40\%$ at
        L1. When the fixed latency dominates, the processor must wait for it anyway and
        neither technique can help much. Here the two levels start out comparable, so
        this is not what separates them.
  \item \emph{The surcharge critical-word first itself costs.} Reordering the transfer
        is not free: $t_c-t_1$ is $20$ CCs at L2 but $9$ CCs at L1. Relative to the
        block service time that is $20/160=12.5\%$ at L2 against $9/30=30\%$ at L1. This
        is the decisive factor: it consumes $12.5$ of L2's $43.75$ available points but
        $30$ of L1's $40$, leaving $31.25\%$ and $10\%$ respectively. The shorter a
        level's total miss time, the harder a fixed reordering overhead is to justify.
  \item \emph{Shape of the reference distribution.} Early restart only pays off to the
        extent that the requested word arrives early. L2's $30\%$ weight on the first
        word helps it: under a uniform $1/8$ the L2 early-restart average would be
        $125$ CCs, a saving of $21.875\%$ rather than $26.25\%$. Critical-word first is
        immune to this, since its cost is flat in the word position --- which makes it
        the more attractive of the two when the distribution is unfavourable or unknown.
  \item \emph{Miss frequency (a full-AMAT caveat).} The comparison above is
        \emph{per miss}, which is what this problem measures. A complete average-memory-
        access-time analysis would also weight each level's saving by how often that
        level misses, and L1 misses are far more frequent than L2 misses. That is the
        one factor that could push the balance back towards L1; with the numbers given
        here, though, L1's saving is only 3 CCs against L2's 50 CCs, so it would take an
        L1:L2 miss ratio above roughly $16{:}1$ merely to draw level.
\end{enumerate}
